\subsubsection{Stationary Laws of Random Function Iterations}

\citet[Theorem~1]{diaconis1999iterated} represent a Markov chain as
$\mathbf{s}_{t+1}=\mathcal F_{\xi_{t+1}}(\mathbf{s}_t)$ with i.i.d.
random Lipschitz maps. On a complete separable metric space, the theorem
requires a finite expected Lipschitz factor, a finite expected displacement
from some reference state, and contraction on average:
\begin{equation}
 \mathbb E[\log K_\xi]<0.
\end{equation}
Here $K_\xi$ bounds the amplification of a state difference by the sampled
map. The conclusion is a unique stationary distribution and geometric
convergence in the Prokhorov distance, a distance between probability laws.
The benchmark's uniform bound $K_\xi\leq\gamma<1$, together with the
bounded invariant domain, supplies a stronger sufficient construction.
The result explains why persistent innovations can sustain fluctuations
while the distribution forgets the initial state. Forward simulations
approach a stationary law; coupling two starting histories with the same
noise gives a separate comparison of their realized paths.
